\chapter{\IfLanguageName{dutch}{Inleiding}{Introduction}}
\label{ch:inleiding}

De initiële projectopstart van een nieuwe frontend-applicatie vraagt vandaag nog veel repetitief werk. Voor digital agencies is niet alleen de \textit{time-to-market} van het eindproduct belangrijk, maar ook de \textit{time-to-demo} in presales- en projectcontext. Bij Ballistix Digital worden B2B-webapplicaties ontwikkeld volgens een gestandaardiseerde architectuur, gebaseerd op een interne React-componentenbibliotheek en vaste code-stijlen. Hoewel veel bouwstenen, zoals navigatiestructuren, formulieren en tabellen, in principe herbruikbaar zijn, verloopt de opstart vandaag nog grotendeels manueel. Dat repetitieve karakter leidt tot onnodig tijdsverlies en vergroot het risico op inconsistenties tussen projecten. Bovendien vereisen handmatige fouten of afwijkingen van de standaard achteraf vaak nog extra correctiewerk.

\section{\IfLanguageName{dutch}{Probleemstelling}{Problem Statement}}
\label{sec:probleemstelling}
Binnen de softwaresector wordt steeds vaker gekeken naar automatisering en AI-ondersteuning als mogelijke oplossingsrichting om dit soort repetitieve taken te vereenvoudigen. Tegelijk blijft het een uitdaging om dergelijke aanpakken af te stemmen op een specifieke bedrijfscontext. Generieke tools en modellen kunnen snel output genereren, maar houden vaak onvoldoende rekening met interne componentenbibliotheken, architecturale afspraken en huisstijlrichtlijnen. Dit resulteert in output die functioneel mogelijk bruikbaar is, maar die nog aanzienlijk extra correctiewerk vereist om volledig binnen de bestaande standaard te passen. 

In die context verkent deze bachelorproef in welke mate context-gedreven aanpakken, eventueel ondersteund door technologieën zoals het Model Context Protocol (MCP), de bedrijfscontext beter kunnen integreren. Deze bijkomende correcties vergroten immers de totale inspanning en beperken de verwachte winst van de automatisering.

\section{\IfLanguageName{dutch}{Onderzoeksvraag}{Research question}}
\label{sec:onderzoeksvraag}
De centrale onderzoeksvraag van deze bachelorproef vermijdt de veronderstelling dat één specifieke technologie de enige oplossing is, en focust in plaats daarvan op het resultaat:
\textit{Welke aanpak biedt binnen een specifieke agency-huisstijl de beste balans tussen snelheid, kwaliteit en consistentie bij de initiële projectopstart van frontend-applicaties?}

De volgende deelvragen ondersteunen het onderzoek:
\begin{enumerate}
    \item \textbf{Deelvraag 1 (Analyse):} Welke stappen in de huidige projectopstart bij Ballistix Digital zijn het meest repetitief en tijdrovend?
    \item \textbf{Deelvraag 2 (Vergelijking):} In welke mate verschillen generieke en context-geoptimaliseerde aanpakken in hun vermogen om consistente output te genereren?
    \item \textbf{Deelvraag 3 (Input-validatie):} Onder welke voorwaarden heeft de meegeleverde context invloed op de bruikbaarheid van de gegenereerde projectonderdelen?
    \item \textbf{Deelvraag 4 (Evaluatie):} Welke aanpak vereist de minste correcties en resulteert in de beste balans qua totale doorlooptijd?
\end{enumerate}

\section{\IfLanguageName{dutch}{Onderzoeksdoelstelling}{Research objective}}
\label{sec:onderzoeksdoelstelling}
Het doel van dit onderzoek is om verschillende aanpakken voor de initiële projectopstart te analyseren en met elkaar te vergelijken binnen de context van Ballistix Digital. Op basis daarvan wordt een \textbf{Proof-of-Concept (PoC)} ontwikkeld om te evalueren in welke mate bedrijfsspecifieke context de kwaliteit, consistentie en efficiëntie van de output kan verbeteren. Daarbij wordt onderzocht of een context-gedreven aanpak daadwerkelijk resulteert in minder correctiewerk en een betere aansluiting op de bestaande architectuur.

\section{\IfLanguageName{dutch}{Opzet van deze bachelorproef}{Structure of this bachelor thesis}}
\label{sec:opzet-bachelorproef}
De rest van deze bachelorproef is als volgt opgebouwd:

In Hoofdstuk~\ref{ch:stand-van-zaken} wordt de stand van zaken binnen het domein besproken, met focus op de balans tussen snelheid en kwaliteit, en relevante benaderingen voor contextvoorziening.

In Hoofdstuk~\ref{ch:methodologie} wordt de methodologie toegelicht, inclusief de nulmeting en de opzet van de experimenten.

In Hoofdstuk~\ref{ch:poc} wordt de technische realisatie van de PoC beschreven.

In Hoofdstuk~\ref{ch:conclusie} worden de resultaten geëvalueerd en de onderzoeksvragen beantwoord.